\documentclass[../../../../main.tex]{subfiles}
\begin{document}

\begin{flushright}
\footnotesize\textit{【自動文字起こし・要確認】}
\end{flushright}

[解] (1)
\[
f_n(x) = \cos\left(x + \frac{a_{n+1}+a_n}{2}\right) \sin\frac{a_{n+1}-a_n}{2} = \frac{1}{2} \left[ \sin(x+a_{n+1}) - \sin(x+a_n) \right]
\]
だから、$h_n(x) = \sum_{k=1}^n f_k(x)$ とおくと
\[
\begin{aligned}
h_n(x) &= \frac{1}{2} \left[ \sin(x+a_{n+1}) - \sin(x+a_1) \right] \quad \cdots \star \\
&= \frac{1}{2} \left[ \sin x \cos a_{n+1} + \sin a_{n+1} \cos x - \sin(x+a_1) \right]
\end{aligned}
\]
だから、$h_n(x)$ が $n \to \infty$ で収束する条件は
\[
\sin a_n, \ \cos a_n \text{ が収束すること。}
\]

(ロ) (イ)がみたされる時 $\sin a_n \to \alpha, \ \cos a_n \to \beta \ (n \to \infty)$ とおくと
\[
F(x) = \frac{1}{2} \left( \beta \sin x + \alpha \cos x - \sin(x+a_1) \right)
\]
より
\[
A = \int_0^{\pi/2} F(x) dx = \frac{1}{2} \left[ \left[ -\beta \cos x + \alpha \sin x \right]_0^{\pi/2} - \left[ -\cos(x+a_1) \right]_0^{\pi/2} \right]
\]
とする。

一方、
\[
\begin{aligned}
B &= \sum_{k=1}^n \left( \int_0^{\pi/2} f_k(x) dx \right) = \frac{1}{2} \sum_{k=1}^n \left( \left[ -\cos(x+a_{k+1}) \right]_0^{\pi/2} + \left[ \cos(x+a_k) \right]_0^{\pi/2} \right) \\
&= -\frac{1}{2} \left[ \left[ \cos(x+a_1) \right]_0^{\pi/2} - \left[ \cos(x+a_{n+1}) \right]_0^{\pi/2} \right] \\
&\to \frac{1}{2} [ P - Q ]
\end{aligned}
\]
だから $A = B$ であり、2つの値は等しい。

\noindent
\textbf{[本時のミス]}\\
・三角は、位相がつれた ずれていればOK!

\end{document}
